\documentclass[12pt,a4paper,twoside,openright,english]{book}
%**************************************************************
% Packages
%**************************************************************
\usepackage[utf8]{inputenc}
\usepackage[T1]{fontenc}
%\usepackage[latin1]{inputenc}
\usepackage[italian]{babel}
\usepackage{graphicx}
\usepackage{float}
\usepackage[toc, acronym]{glossaries}
\usepackage[nottoc]{tocbibind}
\usepackage{transparent}
\usepackage{hyperref}
\usepackage{color}
\usepackage{epstopdf}
\usepackage{subcaption}
\usepackage{wrapfig}
\usepackage{tikz}
%\usepackage{tikz-uml}
%\usepackage{fancyhdr}
\usepackage{rotating}
\usepackage{listings,lstautogobble}
\usepackage{glossaries}
\usepackage{enumitem}



\definecolor{bluekeywords}{rgb}{0,0,1}
\definecolor{greencomments}{rgb}{0,0.5,0}
\definecolor{redstrings}{rgb}{0.64,0.08,0.08}
\definecolor{xmlcomments}{rgb}{0.5,0.5,0.5}
\definecolor{types}{rgb}{0.17,0.57,0.68}



\colorlet{punct}{red!60!black}
\definecolor{background}{HTML}{EEEEEE}
\definecolor{delim}{RGB}{20,105,176}
\colorlet{numb}{magenta!60!black}


\lstdefinelanguage{json}{
	basicstyle=\normalfont\ttfamily,
	numbers=left,
	numberstyle=\scriptsize,
	stepnumber=1,
	numbersep=8pt,
	showstringspaces=false,
	breaklines=true,
	frame=lines,
	backgroundcolor=\color{background},
	literate=
	*{0}{{{\color{numb}0}}}{1}
	{1}{{{\color{numb}1}}}{1}
	{2}{{{\color{numb}2}}}{1}
	{3}{{{\color{numb}3}}}{1}
	{4}{{{\color{numb}4}}}{1}
	{5}{{{\color{numb}5}}}{1}
	{6}{{{\color{numb}6}}}{1}
	{7}{{{\color{numb}7}}}{1}
	{8}{{{\color{numb}8}}}{1}
	{9}{{{\color{numb}9}}}{1}
	{:}{{{\color{punct}{:}}}}{1}
	{,}{{{\color{punct}{,}}}}{1}
	{\{}{{{\color{delim}{\{}}}}{1}
	{\}}{{{\color{delim}{\}}}}}{1}
	{[}{{{\color{delim}{[}}}}{1}
	{]}{{{\color{delim}{]}}}}{1},
}


\lstset{language=[Sharp]C,
	%captionpos=b,
	%numbers=left, %Nummerierung
	%numberstyle=\tiny, % kleine Zeilennummern
	%frame=lines, % Oberhalb und unterhalb des Listings ist eine Linie
	aboveskip=3mm,
	belowskip=3mm,
	showspaces=false,
	showtabs=false,
	breaklines=true,
	showstringspaces=false,
	breakatwhitespace=true,
	escapeinside={(*@}{@*)},
	commentstyle=\color{greencomments},
	morekeywords={partial, var, value, get, set, Public, HttpResponseMessage, async, HttpStatusCode, private, string, List<T>,OrganizationData,Organization},
	keywordstyle=\color{bluekeywords},
	stringstyle=\color{redstrings},
	basicstyle={\ttfamily},
	columns=flexible,	
	breaklines=true,
	tabsize=4,
	autogobble=true
}

\usepackage{setspace}
%\linespread{0.5}

%\pagestyle{fancy}


\usetikzlibrary{calc,trees,positioning,arrows,chains,shapes.geometric,decorations.pathreplacing,decorations.pathmorphing,shapes,matrix,shapes.symbols,fit}



%**************************************************************
% Setup
%**************************************************************

% Line spaceing
\renewcommand{\baselinestretch}{1.5}

% Links color
\hypersetup {
	colorlinks=true,
	linkcolor=black,
	urlcolor=blue,
	citecolor=black
}

% Macros
%\newcommand{\bff}{\acrshort{bff}}

% Diagrams
\tikzset{
	>=stealth',
	punktchain/.style={
		rectangle, 
		rounded corners, 
		fill=yellow!20,
		draw=black, very thick,
		text width=10em, 
		minimum height=3em, 
		text centered, 
		on chain},
	line/.style={draw, thick, <-},
	element/.style={
		tape,
		top color=white,
		bottom color=blue!50!black!60!,
		minimum width=8em,
		draw=blue!40!black!90, very thick,
		text width=10em, 
		minimum height=3.5em, 
		text centered, 
		on chain},
	every join/.style={->, thick,shorten >=1pt},
	decoration={brace},
	tuborg/.style={decorate},
	tubnode/.style={midway, right=2pt},
}

%**************************************************************
% Cover
%**************************************************************
\newcommand{\myName}{Andrea \textsc{Venier}}
\newcommand{\myTitle}{Integrazione tra applicazioni web mediante microservizi RESTful}
\newcommand{\mySubTitle}{Ampliamento delle funzionalità\\ di un’applicazione di Project Management}
\newcommand{\myDegree}{Corso di Laurea in Informatica}
\newcommand{\myUni}{Università di Padova}
\newcommand{\myDepartment}{Dipartimento di Matematica}
\newcommand{\myProf}{Gilberto \textsc{Filè}}
\newcommand{\myLocation}{Padova}
\newcommand{\myAA}{2015-2016}
\newcommand{\myTime}{Dicembre 2016}
\newcommand{\myCompany}{Athesys s.r.l}
\newcommand{\glo}[1]{#1{\ped{G}}}



\makeglossary
%\makenoidxglossaries
%**************************************************************
% Glossary definition
%**************************************************************
\newglossaryentry{package} {
	name=package,
	description={ in informatica è, un raggruppamento di classi, metodi programmi, librerie e procedure che sono logicamente collegate tra di loro
	},
	plural=packages
}
\newglossaryentry{task} {
	name=task,
	description={ è un compito secondo la definizione dello standard ISO/IEC 12207
	},
	plural=tasks
}
\newglossaryentry{Gantt} {
	name=Gantt,
	description={  è un diagramma di supporto alla gestione dei progetti
	}
}
\newglossaryentry{custom} {
	name=custom,
	description={ è un attributo con cui si indica un manufatto, un dispositivo, o un componente, progettato e realizzato su misura in base alle necessità dell'acquirente o della funzione specifica che è destinato ad assolvere
	}
}
\newglossaryentry{APIg} {
	name={API},
	description={ è l'acronimo per Application Programming Interface (in italiano interfaccia di programmazione di un'applicazione), in informatica, indica ogni insieme di procedure disponibili al programmatore, di solito raggruppate a formare un set di strumenti specifici per l'espletamento di un determinato compito all'interno di un certo programma. Spesso con tale termine si intendono le librerie software disponibili in un certo linguaggio di programmazione
	}
}

\newglossaryentry{provider} {
	name=provider,
	description={ è un'espressione inglese che indica imprese che forniscono servizi di vario tipo,  italiano si è soliti tradurre "service provider" letteralmente, chiamando l'impresa "fornitore di servizi"
	},
	plural=providers
}
\newglossaryentry{token} {
	name=token,
	description={ con questo termine si fa riferimento ai Json Web Token (o JWT) che sono stringhe di caratteri alfanumerici che tipicamente incapsulano le credenziali di sicurezza per una sessione di login, identificando l'utente della stessa. Questi token sono segnati da una chiave del server grazie la quale è possibile verificarne la veridicità
	}
}
\newglossaryentry{header} {
	name=header,
	description={ o intestazione, è una parte del pacchetto ,inviato nelle richieste http, che contiene le informazioni di controllo necessarie al funzionamento della rete cioè le informazioni di protocollo aggiunte di strato in strato
	},
	plural=headers
}
\newglossaryentry{endpoint} {
	name=endpoint,
	description={ è l'url attraverso il quale un'applicazione client può accedere ad un servizio web. Lo stesso servizio web di solito ha endpoint multipli
	},
	plural=endpoints
}
\newglossaryentry{JSONg} {
	name={JSON},
	description={  acronimo di Javascript Object Notation, è un formato adatto all'interscambio di dati fra applicazioni client-server
	}
}

\newglossaryentry{REST-based} {
	name=REST-based,
	description={ fa riferimento ad architetture o chiamate  basate sul protocollo REST (vedi sez. \ref{rest})
	}
}
\newglossaryentry{JOIN} {
	name=JOIN,
	description={  è una clausola del linguaggio SQL che serve a combinare (unire) le tuple di due o più relazioni di un database tramite l'operazione di congiunzione (od unione) dell'algebra relazionale
	}
}
\newglossaryentry{SQL} {
	name=SQL,
	description={ è un linguaggio standardizzato per database basati sul modello relazionale
	}
}
\newglossaryentry{query} {
	name=query,
	description={ indica l'interrogazione da parte di un utente di un database, strutturato tipicamente secondo il modello relazionale, per compiere determinate operazioni sui dati
	},
	plural=queries
}
\newglossaryentry{framework} {
	name=framework,
	description={  architettura o struttura di supporto su cui un programma può essere creato. In genere è composto da una serie di librerie e strumenti di sviluppo
	},
	plural=frameworks
}
\newglossaryentry{virtual machine} {
	name=virtual machine,
	description={ o macchina virtuale (VM), viene indicato un software che, attraverso un processo di virtualizzazione, crea un ambiente virtuale che emula tipicamente il comportamento di una macchina fisica grazie all'assegnazione di risorse hardware
	},
	plural=virtual machines
}
\newglossaryentry{Design Pattern} {
	name=Design Pattern,
	description={ nell'ambito dell'ingegneria del software, un design pattern, è un concetto che può essere definito "una soluzione progettuale generale ad un problema ricorrente". Si tratta di una descrizione o modello logico da applicare per la risoluzione di un problema che può presentarsi in diverse situazioni durante le fasi di progettazione e sviluppo del software, ancor prima della definizione dell'algoritmo risolutivo della parte computazionale
	},
	plural=Design Patterns
}
\newglossaryentry{stateless} {
	name=stateless,
	description={ è un protocollo di comunicazione che tratta ogni richiesta come una transazione indipendente, scollegata da qualsiasi precedente richiesta, rendendo la comunicazione composta da coppie indipendenti di richiesta e risposta. Un protocollo stateless inoltre non richiede che il server mantenga le informazioni della sessione per ogni partner di comunicazione per la durata di multiple richieste
	}
}
\newglossaryentry{caching} {
	name=caching,
	description={ è il processo di immagazzinamento dati in una cache, che è una memoria temporanea
	}
}
\newglossaryentry{UML} {
	name=UML,
	description={ è un linguaggio di modellazione e specifica basato sul paradigma object-oriented
	}
}
\newglossaryentry{route} {
	name=route,
	description={ ci si riferisce alla definizione di un URI e a come esso risponde ad una specifica richiesta HTTP di un client
	},
	plural=routes
}
\newglossaryentry{generic} {
	name=generic,
	description={ è un tipo di dato utilizzato negli algoritmi di programmazione che non ha la necessità di essere istanziato immediatamente, bensì può essere specificato in un secondo momento
	},
	plural=generics
}
\newglossaryentry{accessor} {
	name=accessor,
	description={per il linguaggio di programmazione C\# un accessor di una proprietà (o campo dati di una classe) contiene il codice associato con l'operazione di \textit{getting} (lettura) o il \textit{setting} (scrittura) della proprietà stessa. La dichiarazione un un accessor può contenerne uno di tipo \textit{get}, uno di tipo \textit{set} o entrambi
	},
	plural=accessors
}
\newglossaryentry{getter} {
	name=getter,
	description={ è un metodo di una classe utilizzato per leggere il valore di un campo dati della stessa
	},
	plural=getters
}
\newglossaryentry{setter} {
	name=setter,
	description={ è un metodo di una classe utilizzato per scrivere il valore di un campo dati della stessa
	},
	plural=setters
}

\newglossaryentry{URLg} {
	name={URL},
	description={ o Uniform Resource Locator è una sequenza di caratteri che identifica univocamente l'indirizzo di una risorsa in Internet, tipicamente presente su un host server
	}
}
\newglossaryentry{Debugger} {
	name=Debugger,
	description={ è un programma/software specificatamente progettato per l'analisi e l'eliminazione dei bug (debugging), ovvero errori di programmazione interni al codice di altri programmi. E' spesso compreso all'interno di un ambiente integrato di sviluppo (IDE)
	},
	plural=Debuggers
}
\newglossaryentry{logging} {
	name=logging,
	description={ è la registrazione sequenziale e cronologica delle operazioni effettuate, da un utente, un amministratore o automatizzate, man mano che vengono eseguite dal sistema o applicazione
	}
}

\newglossaryentry{asset} {
	name=asset,
	description={ è una risorsa di un azienda
	}
}


\newglossaryentry{URIg} {
	name={URI},
	description={ o Uniform Resource Identifier, in informatica, si riferisce a una stringa che identifica univocamente una risorsa generica che può essere un indirizzo Web, un documento, un'immagine, un file, un servizio, un indirizzo di posta elettronica, ecc
	}
}

\newglossaryentry{cloud-based} {
	name=cloud-based,
	description={ si indica un paradigma basato, sul Cloud, di erogazione di risorse informatiche, come l'archiviazione, l'elaborazione o la trasmissione di dati, caratterizzato dalla disponibilità on-demand attraverso Internet a partire da un insieme di risorse preesistenti e configurabili
	}
}

%**************************************************************
% Acronyms definitions
%**************************************************************
\newacronym{CRM}{CRM}{Customer relationship management}

%\newacronym{JSON}{JSON}{JavaScript Object Notation}

\newacronym{http}{HTTP}{Hypertext Transfer Protocol}

\newacronym{DAO}{DAO}{Data Access Object}

\newacronym{DTO}{DTO}{Data Transfer Object}

%\newacronym{URI}{URI}{Uniform Resource Identifier}





%%% define the acronym and use the see= option
\newglossaryentry{URI}{type=\acronymtype, name={URI}, description={Uniform Resource Identifier}, first={Uniform Resource Identifier (URI)\glsadd{URIg}}, see=[Glossary:]{URIg}}

\newglossaryentry{JSON}{type=\acronymtype, name={JSON}, description={JavaScript Object Notation}, first={JavaScript Object Notation (JSON)\glsadd{JSONg}}, see=[Glossary:]{JSONg}}

\newglossaryentry{API}{type=\acronymtype, name={API}, description={Application Programming Interface}, first={Application Programming Interface (API)\glsadd{APIg}}, see=[Glossary:]{APIg}}

\newglossaryentry{URL}{type=\acronymtype, name={URL}, description={Uniform Resource Locator}, first={Uniform Resource Locator (URL)\glsadd{URLg}}, see=[Glossary:]{URLg}}

\title{\myTitle}
\author{\myName}
\date{\today}

\begin{document}

%**************************************************************
% Cover
%**************************************************************
\frontmatter
\begin{titlepage}
	\begin{center}
		\vbox to40pt{\vbox to\textheight{\vfill {\transparent{0.2}\includegraphics[width=10cm]{logo-unipd.png}} \vfill}}
	\end{center}
	\begin{minipage}{.20\textwidth}
		\includegraphics[height=2cm]{logo-department.png}
	\end{minipage}
	\hspace{40pt}
	\begin{minipage}{.80\textwidth}
		\begin{center}
			\begin{LARGE}
				\textsc{\textbf{\myUni}}\\
			\end{LARGE}
			\line(1,0){220}\\
			\begin{Large}
				\textsc{\myDepartment}\\
			\end{Large}
		\end{center}
	\end{minipage}
	\begin{center}
		\vfill
		\large{\textsc{\myDegree}}\\			
		\LARGE{\textsc{\textbf{\myTitle}}}\\
		\large{\mySubTitle}\\
		\vfill
		\begin{minipage}{0.4\textwidth}
			\begin{flushleft} \large
				\emph{Autore}\\
				\myName
			\end{flushleft}
		\end{minipage}
		\hfill
		\begin{minipage}{0.4\textwidth}
			\begin{flushright} \large
				\emph{Relatore} \\
				\myProf
			\end{flushright}
		\end{minipage}
		\vfill
		\line(1, 0){338} \\
		\begin{normalsize}
			\textsc{Anno Accademico \myAA}
		\end{normalsize}
	\end{center}
\end{titlepage}

%**************************************************************
% Colophon
%**************************************************************
\thispagestyle{empty}
\hfill
\vfill
\noindent \myName\ \textit{\myTitle,} \myDegree, Copyright \textcopyright\ \myTime.

\cleardoublepage

%**************************************************************
% Abstract
%**************************************************************
\chapter*{Abstract}\label{abstract}

%**************************************************************
% Keywords
%**************************************************************
\chapter*{Keywords}\label{keywords}
\paragraph*{CRM}
sta per \textit{Customer relationship management} e rappresenta un approccio per gestire le interazioni di un azienda con i suoi clienti, attuali e potenziali.\\
Attraverso questa metodologia si punta a raccogliere e analizzare i dati relativi al cliente (email, sito aziendale, numeri telefonici, offerte precedenti) con il fine di incrementare le vendite grazie alla fidelizzazione dello stesso mediante offerte mirate.\\
Seguendo questi concetti sono nati, e si sono affermati, diversi software CRM e tra i maggiori esponenti possiamo trovare: \textbf{Salesforce} e \textbf{Microsoft Dynamics}.\\
L'azienda \large{\myCompany{}}, basandosi sull'importanza e sulla diffusione dei suddetti software ha deciso di integrarli nella propria applicazione web di \textit{Project Management}.
 
\paragraph*{Project Management Software}
La disciplina di \textit{Project Management} si occupa di sovraintendere la pianificazione, l'organizzazione e l'implementazione di un progetto. Dove con progetto si intende un insieme ordinato di compiti da svolgere, caratterizzati da: vincoli di tempo (date di inizio e fine), risorse utilizzabili e risultati attesi.\\
Questa disciplina è fondamentale per controllare i costi, l'utilizzo delle risorse aziendali e per poter assicurare che il prodotto finale di questi processi soddisfi a pieno le esigenze del cliente.\\
In questo contesto \large{\myCompany{}} ha deciso di sviluppare \textbf{ADProject}, una soluzione software pensata ad-hoc per i propri clienti
l'obbiettivo di facilitare e automatizzare il suddetto insieme di attività.


%**************************************************************
% Table of contents, list of figures and list of tables
%**************************************************************
\tableofcontents
\listoffigures
\listoftables

\mainmatter

%**************************************************************
% Introduzione
%**************************************************************
\chapter{Introduzione}\label{introduzione}
\section{Il progetto}
%**************************************************************
% Tecnologie e Strumenti
%**************************************************************
\chapter{Tecnologie e Strumenti}\label{tecnologie}
\section{Frameworks e linguaggi di programmazione}
In questa sezione verranno descritti brevemente i linguaggi di programmazione usati nella scrittura di \textbf{ADProject} e di \textbf{ADCrm}.
Non tutti i linguaggi che verranno presentati in seguito sono stati usati nell'ambito del progetto di stage (in quanto limitato alla realizzazione del servizio \large{ADCrm}), tuttavia la loro contestualizzazione all'interno dei vari moduli sviluppati dall'azienda aiuta a capire meglio il funzionamento degli stessi.

\paragraph{.NET Framework}
è un framework software sviluppato da Microsoft. Include una vasta classe di librerie chiamata \textit{Framework Class Library} (FLC) che permette l'interoperabilità tra i linguaggi di programmazione del framework. I programmi scritti per .NET vengono eseguiti in un ambiente software conosciuto come \textit{Common Language Runtime} (CLR), una \textit{virtual machine}  che fornisce servizi come la gestione delle eccezioni, gestione della memoria, sicurezza. FCL e CLR insieme costituisco il framework .NET.\\
Il software di Project management \textbf{ADProject} e il modulo \textbf{ADCrm} sono sviluppati su questo framework.
%TODO da incrementare questa parte o specificarla meglio

\paragraph{ASP.NET}
 è un framework server-side, sviluppato da Microsoft, per la creazione di applicazioni web che permette ai programmatori di costruire pagine web dinamiche, applicazioni e servizio web.
 ASP.NET viene usato per lo sviluppo di tutte le pagine web dinamiche che fungono da interfacce grafiche per ADProject.

\paragraph{Visual Basic .NET}
è un linguaggio di programmazione orientato agli oggetti sviluppato sul framework .NET.\\
Con questo linguaggio è stato sviluppato ADProject, e sono tuttora in sviluppo numerosi moduli di questo software. La motivazione principale che ha portato i developer assegnati al progetto a scegliere questo linguaggio è stata la loro famigliarità con lo stesso. 

\paragraph{C\#}
è un linguaggio di programmazione orientato agli oggetti sviluppato sul framework .NET.
La sintassi del linguaggio prende spunto da C++ , da Java
e da Visual Basic.
Questo linguaggio sta alla base di tutto lo sviluppo del microservice ADCrm, scelta che si potrebbe pensare poco sensata in quanto ADCrm deve integrarsi con un applicazione web interamente scritta con Visual Basic .NET, tuttavia ci sono delle forti motivazioni che hanno fatto propendere il team di sviluppo per adottare questo linguaggio:
\begin{itemize}
	\item avendo scelto di adottare una struttura basata su microservices per lo sviluppo di ADCrm, ed essendo quindi un modulo fisicamente slegato da ADProject, è possibile scegliere un linguaggio di programmazione differente da VB.NET, pur rimanendo all'interno dei linguaggi offerti dal framework .NET;
	\item i programmatori assegnati allo sviluppo del progetto hanno maggiore famigliarità con questo linguaggio;
	\item Microsoft tendenzialmente offre maggiore supporto per C\#;
	\item C\# ha una vastissima comunity di sviluppatori, persino superiore a quella di VB.NET;
	\item è più facile reperire documentazione e teamplate rispetto a VB.NET;
	\item è più facile trovare nuovi sviluppatori per C\#.	
\end{itemize}

\section{Ambiente di sviluppo}
\paragraph{Visual Studio}
 è un potente l'ambiente di sviluppo multipiattaforma fornito da Microsoft ed è stata la scelta naturale per lo sviluppo di ogni componente di ADProject ed ADCrm.
Visual Studio supporta, oltre ai linguaggi e framework sopracitati: C, C++, F\#, Html e Javascript. 
%TODO da specificare cosa vuol dire multipiattaforma

\chapter{Analisi}\label{analisi}

%**************************************************************
% Casi d'uso
%**************************************************************
\section{Casi d'uso}\label{usecase}
%TODO mancano gli include

\begin{small}

	\hypertarget{UC1}{}
	\subsection{Caso d'uso UC1: Configurazione Aziendale}
	\begin{longtable}{ | p{2.7cm} | p{12cm} |}
		\hline \textbf{Attori} & Utente con privilegi di scrittura\\ 
		\hline \textbf{Descrizione} & Attraverso l’interfaccia di configurazione aziendale, l’utente coi permessi corretti è in grado di visualizzare le impostazioni riguardanti la configurazione aziendale e modificarle\\ 
		\hline \textbf{Scenario Principale} & \begin{enumerate}
			\item L’utente visualizza la lista delle famiglie di prodotti presenti nel listino del CRM  (UC1.1);
			\item L’utente seleziona la tipologia di risorse con cui andare a mappare ciascuna famiglia di prodotti  (UC1.2);
			\item L’utente mappa eventuali famiglie di prodotti non mappate o prodotti specifici di un’offerta che non sono stati mappati  (UC1.3);
			
		\end{enumerate}
		\\ 
		\hline \textbf{Postcondizione} & Sono state visualizzate ed eventualmente modificate le associazioni tra le famiglie di prodotti e le tipologie di risorse\\ 
		\hline 
	\end{longtable}
	
	\hypertarget{UC1.2}{}
	\subsection{Caso d'uso UC1.2: Selezione della tipologie di risorsa con cui mappare ciascuna famiglia di prodotti}
	\begin{longtable}{ | p{2.7cm} | p{12cm} |}
		\hline \textbf{Attori} & Utente con privilegi di scrittura\\ 
		\hline \textbf{Descrizione} & Attraverso l’interfaccia di configurazione aziendale, l’utente coi permessi corretti è in grado di specificare con quale tipologia di risorsa (figura professionale, asset o servizio) dovrà essere mappata ciascuna tipologia di prodotto proveniente dal CRM\\ 
		\hline \textbf{Scenario Principale} & \begin{enumerate}
			\item L’utente mappa una famiglia di prodotti con le figure professionali necessarie per un progetto  (UC1.2.1);
			\item L’utente mappa una famiglia di prodotti con gli asset necessarie per un progetto  (UC1.2.2);
			\item L’utente mappa una famiglia di servizi necessari per un progetto  (UC1.2.3);
			\item L’utente mappa una famiglia di prodotti come prodotti venduti (UC1.2.4);
			
		\end{enumerate}
		\\ 
		\hline \textbf{Estensioni} & \begin{enumerate}
			\item L’utente visualizza un messaggio di errore se tenta di cambiare la mappatura di una famiglia dopo che sono già stati importati prodotti relativi a quella famiglia  (UC1.4);
			
		\end{enumerate}
		\\ 
		\hline \textbf{Postcondizione} & Per almeno una famiglia di prodotto è stata modificata l’associazione ad una tipologia di risorsa. \\ 
		\hline 
	\end{longtable}
	
	\hypertarget{UC1.3}{}
	\subsection{Caso d'uso UC1.3: Mappatura famiglie di prodotti aggiunte successivamente con prodotti specifici di un offerta}
	\begin{longtable}{ | p{2.7cm} | p{12cm} |}
		\hline \textbf{Attori} & Utente con privilegi di scrittura\\ 
		\hline \textbf{Descrizione} & Attraverso l’interfaccia di configurazione di progetto, l’utente coi permessi corretti è in grado di effettuare una mappatura \textit{just-in-time} di eventuali prodotti legati all’offerta che appartengono ad una famiglia di prodotti non mappata (perché ad esempio è stata aggiunta da poco nel CRM) o che non sono associati ad una famiglia (perché magari sono prodotti associati puntualmente all’offerta\\ 
		\hline \textbf{Scenario Principale} & \begin{enumerate}
			\item L’utente visualizza la lista dei prodotti non associati ad alcuna famiglia del CRM  (UC1.3.1);
			\item L’utente seleziona la categoria di risorsa con cui mappare ciascun prodotto  (UC1.3.2);
			\item L’utente visualizza la lista delle famiglie di prodotti non mappate  (UC1.3.3);
			\item L’utente seleziona la categoria di risorsa con cui mappare ciascuna famiglia  (UC1.3.4);
			
		\end{enumerate}
		\\ 
		\hline \textbf{Postcondizione} & Sono state potenzialmente modificate le mappature dei prodotti\\ 
		\hline 
	\end{longtable}
	
	\hypertarget{UC2}{}
	\subsection{Caso d'uso UC2: Gestione delle anagrafiche aziendali}
	\begin{longtable}{ | p{2.7cm} | p{12cm} |}
		\hline \textbf{Attori} & Utente con privilegi di lettura, Utente con privilegi di scrittura\\ 
		\hline \textbf{Descrizione} & Attraverso l’interfaccia di gestione delle anagrafiche aziendali, l’utente coi permessi corretti è in grado di visualizzare le informazioni relative alle schede di anagrafica aziendali e modificarle\\ 
		\hline \textbf{Scenario Principale} & \begin{enumerate}
			\item L’utente crea un’anagrafica di un’azienda  (UC2.1);
			\item L’utente modifica l’anagrafica di un’azienda già censita  (UC2.3);
			\item L’utente gestisce i contatti relativi ad un’azienda  (UC2.4);
			\item L’utente visualizza la lista delle aziende censite nel sistema  (UC2.5);
			\item L’utente visualizza i dettagli relativi ad una azienda censita nel sistema  (UC2.6);
			
		\end{enumerate}
		\\ 
		\hline \textbf{Postcondizione} & Sono state visualizzate ed eventualmente modificate le anagrafiche aziendali\\ 
		\hline 
	\end{longtable}
	
	\hypertarget{UC2.4}{}
	\subsection{Caso d'uso UC2.4: Gestione dei contatti relativi ad un'azienda}
	\begin{longtable}{ | p{2.7cm} | p{12cm} |}
		\hline \textbf{Attori} & Utente con privilegi di scrittura\\ 
		\hline \textbf{Descrizione} & Attraverso l’interfaccia di gestione delle anagrafiche aziendali, l’utente coi permessi corretti è in grado di visualizzare le informazioni sui contatti di riferimento per le diverse aziende ed eventualmente modificarle\\ 
		\hline \textbf{Scenario Principale} & \begin{enumerate}
			\item L’utente crea un’anagrafica di contatto ;
			\item L’utente modifica l’anagrafica di un contatto già censito ;
			\item L’utente visualizza la lista dei contatti di riferimento per una specifica azienda ;
			\item L’utente visualizza i dettagli relativi ad un contatto ;
			
		\end{enumerate}
		\\ 
		\hline \textbf{Postcondizione} & Sono state visualizzate ed eventualmente modificate le informazioni relative ai contatti di riferimento per un’azienda\\ 
		\hline 
	\end{longtable}
	
	\hypertarget{UC3}{}
	\subsection{Caso d'uso UC3: Gestione dei progetti}
	\begin{longtable}{ | p{2.7cm} | p{12cm} |}
		\hline \textbf{Attori} & Utente con privilegi di lettura, Utente con privilegi di scrittura\\ 
		\hline \textbf{Descrizione} & Attraverso l’interfaccia di configurazione di progetto, l’utente coi permessi corretti è in grado di modificare le informazioni di progetto \\ 
		\hline \textbf{Scenario Principale} & \begin{enumerate}
			\item L’utente associa un commerciale ad un progetto  (UC3.1);
			\item L’utente associa una delle aziende clienti censite al progetto  (UC3.2);
			\item L’utente associa una o più offerte libere al progetto selezionato  (UC3.3);
			\item L’utente gestisce il contatto principale di riferimento per un progetto  (UC3.4);
			\item L’utente visualizza la lista delle offerte libere (non associate ad alcun progetto) relative al cliente selezionato  (UC3.6);
			\item L’utente visualizza i dettagli delle offerte associate al progetto  (UC3.7);
			\item L’utente visualizza la lista di tutte le offerte libere  (UC3.8);
			\item L’utente può creare un progetto a partire da un’offerta  (UC3.5);
			
		\end{enumerate}
		\\ 
		\hline \textbf{Postcondizione} & Sono modificate le informazioni relative ad un progetto. \\ 
		\hline 
	\end{longtable}
	
	\hypertarget{UC3.4}{}
	\subsection{Caso d'uso UC3.4: Gestione del contatto principale del progetto}
	\begin{longtable}{ | p{2.7cm} | p{12cm} |}
		\hline \textbf{Attori} & Utente con privilegi di scrittura\\ 
		\hline \textbf{Descrizione} & Attraverso l’interfaccia di configurazione di progetto, l’utente coi permessi corretti è in grado di modificare le informazioni relative al contatto principale di riferimento per un progetto (il PM lato cliente)\\ 
		\hline \textbf{Scenario Principale} & \begin{enumerate}
			\item L’utente seleziona uno dei contatti associato al cliente selezionato per il progetto  (UC3.4.1);
			\item L’utente inserisce un nuovo contatto  (UC3.4.2);
			\item L’utente modifica i dettagli relativi al contatto di riferimento per il progetto  (UC3.4.3);
			\item L’utente visualizza i dettagli relativi al contatto inserito  (UC3.4.4);
			
		\end{enumerate}
		\\ 
		\hline \textbf{Estensioni} & \begin{enumerate}
			\item Viene visualizzato un messaggio d'errore se si tenta di aggiungere un contatto principale prima che sia associato un'azienda cliente al progetto (UC3.12);
			
		\end{enumerate}
		\\ 
		\hline \textbf{Postcondizione} & Sono state modificate le informazioni relative al contatto principale di riferimento per un progetto \\ 
		\hline 
	\end{longtable}
	
	\hypertarget{UC3.5}{}
	\subsection{Caso d'uso UC3.5: Creazione nuovo progetto da offerta}
	\begin{longtable}{ | p{2.7cm} | p{12cm} |}
		\hline \textbf{Attori} & Utente con privilegi di scrittura\\ 
		\hline \textbf{Descrizione} & L'utente crea un progetto a partire dai dati di una specifica offerta, auto-completando quindi l'associazione tra progetto con account, offerta e contatto\\ 
		\hline \textbf{Scenario Principale} & \begin{enumerate}
			\item L’utente inserisce i dati di base del progetto  (UC3.5.1);
			\item L’utente selezione il cliente su cui filtrare le offerte  (UC3.5.2);
			\item L’utente seleziona l’offerta da cui creare il nuovo progetto  (UC3.5.3);
			
		\end{enumerate}
		\\ 
		\hline \textbf{Postcondizione} & Viene creato un nuovo progetto\\ 
		\hline 
	\end{longtable}
\end{small}





%**************************************************************
% Requisiti
%**************************************************************
\section{Requisiti}\label{requisiti}


%**************************************************************
% Progettazione
%**************************************************************
\chapter{Progettazione}\label{progettazione}

\section{Design pattern}\label{design_pattern}
\section{Progettazione architetturale}\label{progettazione_architetturale}
\section{Progettazione di dettaglio}\label{progettazione_dettaglio}

\chapter{Verifica e validazione}\label{verifica_validazione}

\chapter{Conclusioni}\label{conclusioni}

%**************************************************************
% Glossary and bibliography
%**************************************************************
\printglossaries
\bibliographystyle{unsrt}
%\bibliography{thesis}


%**************************************************************
% Ringraziamenti
%**************************************************************

%\chapter*{Thanksgivings}
%rr 

%\noindent\textit{\myLocation, \myTime}
%\hfill \myName

\end{document}